\\documentclass[11pt]{article}
\\usepackage{fontspec}
\\usepackage{xeCJK}
\\setmainfont{Times New Roman}
\\setCJKmainfont{Noto Serif CJK SC}
\\usepackage{geometry}
\\geometry{a4paper,margin=1in}
\\usepackage{setspace}
\\onehalfspacing
\\usepackage{microtype}
\\usepackage{hyperref}
\usepackage[backend=biber,style=numeric,sorting=none]{biblatex}
\addbibresource{refs.bib}
\\begin{document}
\title{fb24c320-f14e-4d5d-811e-f1f4204fd0c3-超高压阀门流动特性与响应规律仿真}\maketitle\n
\par
\par
\par
超高压阀门流动特性与响应规律仿真\par

\par
娄宁\par

2026 年 8 月 27  日\par

\par
\par
\textbf{摘要}\par

针对 320MPa 级超高压氮气管路系统，本文采用计算流体力学方法开展了多级节流阀门的瞬态流动特性数值模拟研究。建立了包含高压储罐、三级支管与节流阀芯的三维计算域， 采用 Peng-Robinson 真实气体模型描述超高压氮气物性， 对比了压力基隐式求解器与密度基显式 Roe-FDS 求解器的计算稳定性，分析了不同湍流模型对强瞬态流动的适配性。研究成功捕捉了阀门动作前的稳态流动分布与两级阀门动作过程的瞬态响应规律，得到了压力、速度与质量流量的变化特征。针对强瞬态阶段湍流方程数值发散的问题，从 CFL 稳定性条件、流场梯度与湍流模型特性三个层面揭示了发散机制。研究结果可为超高压阀门系统的设计与响应特性分析提供参考。\par

\textbf{关键词：}超高压氮气；多级节流阀；瞬态流动；数值模拟；计算稳定性\par

\par
\par
\par
\textbf{1}\textbf{    }\textbf{引言}\par

超高压气体管路系统广泛应用于航空航天、特种装备、油气储运与工业储能领域，其中节流阀门是控制系统压力与流量的核心部件 [\textit{,}]。阀门开关过程中的瞬态流动会诱发压力波振荡、流量冲击与气动噪声， 直接影响系统的安全性与控制精度。由于超高压工况下实验测试难度大、成本高， 且存在安全风险， 数值模拟成为揭示阀门瞬态流动规律的重要手段 [\textit{,}]。\par

国内外学者针对高压阀门内部流动开展了大量研究。李华贵等 [\cite{3b7abd61_5375_480a_b603_09b3c731c43f_01_DN89切断阀压损稳态仿真阶段性报告}]  针对超高压大口径气动截止阀开展了结构优化研究， 通过压力平衡式阀瓣设计降低了关闭阻力， 实现了0.02s 的超快开启，验证了超高压阀门在特种装备领域的应用价值。梁连金 [\cite{88d32ae0_631e_44d3_8492_0fcfdf192585_超高压阀门流动特性与响应规律仿真}]  针对高压大口径提升杆球阀， 提出了浮动阀座与容腔缓冲结构， 显著提升了双向密封性能与压力稳定性。工程实践中， 上海电气 [\cite{99fa1095_b0df_4b4c_be5d_f90ea8582a3c_超高压阀门流动特性与响应规律仿真}]  针对超高压氮气管路系统开展了系列仿真研究， 采用Peng-Robinson 真实气体模型完成了 DN89、DN65 等多规格阀门的稳态压损计算， 为工程设计提供了数据支撑 [ \textit{—}]。\par

\par
在数值方法层面，谭宗豪 [\cite{9b9e80d2_e2e9_458b_9dae_18b0303b1c6c_超高压大口径气动快速截止阀的研制_李华贵}]  在管道阀门监测研究中，对比了多种湍流模型与网格方案的计算精度与效率，指出 k-  系列模型在工程流动中具有良好的适用性，但在强梯度工况下易出现数值振荡。原佳妮 [\cite{d9d53741_b3eb_4a59_8e8b_ce34e7f99554_02_DN65阀门稳态压损仿真阶段性报告}]  对高温高压球阀的瞬态特性研究表明， 阀门动作过程中的强压力冲击会诱发流场与结构的耦合振荡，对数值稳定性提出了较高要求。\par

现有研究多针对单级阀门的稳态流动， 对 300MPa 级超高压下多级节流阀的瞬态响应特性与数值稳定性研究较少。本项目以三级节流超高压阀门系统为研究对象， 基于计算流体力学方法开展瞬态流动仿真， 旨在揭示阀门动作过程中的压力传播与流量响应规律，分析强瞬态工况下的数值稳定性问题， 为同类超高压阀门系统的设计与优化提供依据。\par

\par
\textbf{2    }\textbf{几何模型与计算域}\par

研究对象为多级节流式高压阀门管路系统，主要由上游高压储罐、三级分支管路、节流阀芯组件与下游出口管路组成。系统上游入口压力为 320MPa，下游出口压力为220MPa，两级阀门按预设时序依次动作，实现流量的分级调节。该系统与文献 [\cite{f9aa7304_1ffe_4eea_aff7_c3174b482e09_高温高压气动球阀的结构优化及响应特性研究_原佳妮}]   中的汇管段结构类似，但支路更多、压力等级更高。\par

计算域覆盖从储罐出口到下游出口的全部流道，忽略阀体结构的热变形与弹性变形，仅考虑流体域的瞬态流动。计算域采用六面体结构化网格进行离散，总网格量约为 68 万。对节流孔、阀芯壁面等流动梯度较大的区域进行网格加密，最小网格尺寸为0.2mm，以准确捕捉边界层与节流射流特征。该局部加密原则与文献 [\cite{fb24c320_f14e_4d5d_811e_f1f4204fd0c3_超高压阀门流动特性与响应规律仿真}]  提出的曲率尺寸函数方法一致，在保证计算精度的同时控制计算量。网格划分方式同样参考了上海电气[\cite{fd169dda_f9ff_45e3_a3c8_72b02ab15ab5_高压大口径双向密封程控提升杆球阀的研制_梁连金}]  的工程仿真规范，采用结构化网格提升数值稳定性。\par

\par
图 1    系统计算域与初始压力分布云图\par

\par
\textbf{3    }\textbf{数值方法与计算设置}\par

\par
\textbf{3.1    }\textbf{控制方程与物性模型}\par

瞬态可压缩流动满足质量守恒、动量守恒与能量守恒方程， 采用雷诺平均方法对控制方程进行离散求解。由于工作压力高达 320MPa，氮气分子间作用力显著，理想气体状态方程的计算误差随压力升高急剧增大。DN89 与 DN65  阀门的工程仿真结果表明[ \textit{—}] ，当压力超过 30MPa 时，理想气体模型的压损计算偏差可达 15\\% 以上，必须采用真实气体模型。\par

本研究采用 \textbf{Peng}\textbf{-}\textbf{Robinson}\textbf{ }\textbf{真实气体模型}计算密度与压力、温度的耦合关系， 其基本形式为：\par

\par
式中：\textit{P}\textit{ }为压力，\textit{T}\textit{ }为温度，\textit{R}\textit{ }为气体常数，\textit{v}\textit{ }为摩尔体积，\textit{a}、\textit{b}\textit{ }为气体特性参数， \textit{α}\textit{ }(\textit{T})为温度函数。该模型在高压气动系统仿真中已得到广泛的工程验证 [ \textit{,}\textit{,}] ，上海电气 [\cite{fe880629_9ce4_4876_a173_e24a97d827f4_基于数字孪生的管道阀门监测技术研究_谭宗豪}]在全系列阀门仿真中均采用该模型，计算结果与实验偏差在工程允许范围内。\par

\par
\textbf{3.2    }\textbf{湍流模型对比}\par

研究先后测试了三种粘性计算方案，对比其在强瞬态工况下的稳定性与适用性：\par

1.  \textbf{Standard}\textbf{ }\textbf{k}\textbf{-}\textbf{epsilon}\textbf{ }\textbf{模型} \textbf{+}\textbf{ }\textbf{标准壁面函数}：工业通用湍流模型，通过求解湍动能 k 和耗散率  的输运方程封闭雷诺平均方程。谭宗豪 [\cite{fe880629_9ce4_4876_a173_e24a97d827f4_基于数字孪生的管道阀门监测技术研究_谭宗豪}]  在管道阀门流场仿真中指出，该模型对充分发展湍流的计算精度良好， 但在强逆压梯度、流动分离等工况下误差较大。\par

2.  \textbf{Realizable}\textbf{ }\textbf{k}\textbf{-}\textbf{epsilon  }\textbf{模型}：对标准模型的耗散率方程进行了修正，引入了旋转与曲率修正，对强逆压梯度流动的适应性优于标准模型。原佳妮 [\cite{fe880629_9ce4_4876_a173_e24a97d827f4_基于数字孪生的管道阀门监测技术研究_谭宗豪}]  在高压球阀仿真中也采用了该模型， 以提升强冲击工况下的收敛性。计算中同时开启 Production Limiter 限制湍动能的过度产生， 并设置湍流粘度比上限为 10，以提升数值稳定性。\par

3.  \textbf{层流模型}：直接求解层流 N-S 方程，不引入湍流模型假设，从根源上消除了湍流方程的数值发散风险，适用于强瞬态工况下的定性流动规律分析 [\cite{fe880629_9ce4_4876_a173_e24a97d827f4_基于数字孪生的管道阀门监测技术研究_谭宗豪}]。\par

\par
\textbf{3.3    }\textbf{求解器配置}\par

分别采用两类求解器进行试算对比：\par

1.  \textbf{压力基隐式求解器}：采用 SIMPLE 算法实现压力与速度的耦合，通过亚松弛因子控制变量更新幅度，适合低速不可压缩流动。工程上 DN89 、DN65 等阀门的稳态仿真均采用该类求解器 [ \textit{—}] ，计算稳定可靠。\par

\par
2.  \textbf{密度基显式求解器}：直接求解可压缩流动守恒方程， 采用 Roe-FDS 通量离散格式，对激波、压力波等间断流动的捕捉精度更高。对于显式时间推进格式， 其线性稳定性满足 Courant-Friedrichs-Lewy 判据 [\cite{fe880629_9ce4_4876_a173_e24a97d827f4_基于数字孪生的管道阀门监测技术研究_谭宗豪}] ，即 CFL 数不大于 1，时间步长受最小网格尺寸与当地声速的共同约束。\par

\par
\textbf{3.4    }\textbf{边界条件与初始化}\par

•  入口边界：压力入口，总压 320MPa，温度 288K，与文献 [ \textit{,}]  的入口工况一致；\par

•  出口边界：压力出口，静压 220MPa，开启\textbf{防止逆流}选项，避免瞬态过程中回流气体带入非物理扰动，诱发计算发散 [\cite{fe880629_9ce4_4876_a173_e24a97d827f4_基于数字孪生的管道阀门监测技术研究_谭宗豪}] ；\par

•  壁面边界：无滑移绝热壁面， 壁面粗糙度设置为 0.8，与工程仿真标准保持一致 [\cite{fe880629_9ce4_4876_a173_e24a97d827f4_基于数字孪生的管道阀门监测技术研究_谭宗豪}] ；\par

•  初始化：采用\textbf{分区} \textbf{Patch}\textbf{ }\textbf{初始化方法}，上游区域初始化为 320MPa，下游阀门后区域初始化为 220MPa，构建符合物理实际的初始压力场， 避免全场均匀初始化导致的初始时刻巨大压力梯度，有效降低了计算初期的数值振荡 [\cite{fe880629_9ce4_4876_a173_e24a97d827f4_基于数字孪生的管道阀门监测技术研究_谭宗豪}]。\par

\par
\textbf{4    }\textbf{计算稳定性调试过程}\par

超高压强瞬态流动的数值刚性大， 计算易发散， 项目通过多轮参数调整与方案对比逐步优化稳定性。调试思路参考了谭宗豪 [\cite{fe880629_9ce4_4876_a173_e24a97d827f4_基于数字孪生的管道阀门监测技术研究_谭宗豪}]  提出的“亚松弛-限制器-迭代步数”三级优化策略。\par

\par
\textbf{4.1    }\textbf{压力基求解器调试}\par

初始采用压力基求解器配合 Realizable k-epsilon 模型，计算至 0.056s 时湍动能 k出现持续飙升，最终导致数值溢出。针对该问题， 首先将 k、 方程的亚松弛因子由默认的 0.8 降低至 0.3，减小每步迭代的变量更新幅度，抑制数值振荡的累积。优化后计算稳定时间提升至 0.08s。\par

在此基础上开启湍动能产生项限制器（Production Limiter）， 限制壁面附近与强梯度区域的湍动能过度产生， 同时将每时间步最大迭代数从 60 提升至 150，强化单时间步内的收敛性。采取上述措施后，计算稳定时间进一步提升至 0.14s，但在阀门动作的强瞬态阶段，湍流方程仍出现非物理溢出。原佳妮 [\cite{fe880629_9ce4_4876_a173_e24a97d827f4_基于数字孪生的管道阀门监测技术研究_谭宗豪}]  在研究中也发现，强瞬态冲击下，即使采用 Realizable 模型，湍动能方程仍易出现非物理增长。\par

\par
\textbf{4.2    }\textbf{密度基求解器测试}\par

依据标准可压缩流计算方法，切换为密度基显式 Roe-FDS 求解器，配合参考参数（粘度 0.00125 Pa ·s、时间步长 0.001s）进行测试。根据显式格式的 CFL 稳定性判据 [\cite{fe880629_9ce4_4876_a173_e24a97d827f4_基于数字孪生的管道阀门监测技术研究_谭宗豪}] ：\par

\par
\par
式中：\textit{c}\textit{ }为当地声速， ∆\textit{t}\textit{ }为时间步长， ∆\textit{x}\textit{ }为网格特征尺寸。\par

本模型节流口处最小网格尺寸约 0.2mm，超高压下氮气当地声速约为 520m/s，代入公式可得理论稳定时间步长约为 3\textit{.}8 \textit{×}\textit{ }10\textit{—}7  s。采用 0.001s 的参考步长远大于稳定极限约 3 个数量级， 因此计算起步即触发浮点异常。这一现象完全符合显式时间推进格式的基本稳定理论 [\cite{fe880629_9ce4_4876_a173_e24a97d827f4_基于数字孪生的管道阀门监测技术研究_谭宗豪}] ，说明超高压细网格工况下， 显式求解器的计算效率会受到极大限制。\par

\par
\textbf{4.3    }\textbf{层流定性方案}\par

为验证湍流模型是发散根源， 切换为层流模型进行试算。关闭湍流方程后， 计算稳定性显著提升，可完整覆盖阀门动作的瞬态过程，进一步证实了 k-epsilon 湍流模型在强瞬态强梯度工况下的数值不稳定性。层流计算结果可用于分析流动趋势与响应时序，为后续湍流模型优化提供定性参考 [\cite{fe880629_9ce4_4876_a173_e24a97d827f4_基于数字孪生的管道阀门监测技术研究_谭宗豪}]。\par

\par
\par
图 2    典型残差变化曲线（0.14s  内稳定阶段）\par

\par
\par
\textbf{5    }\textbf{计算结果与分析}\par

\par
\textbf{5.1    }\textbf{稳态流动特性（阀门动作前）}\par

在阀门动作前的 0 0.08s，流动逐步建立稳态，流场呈现以下特征：\par

\par
\textbf{5.1.1    }\textbf{压力分布}\par

整体呈现上游高压、下游低压的逐级分布特征。总压降约为 10MPa，其中约 85\\%的压降集中在节流孔位置，主管路与容腔段的压力损失不足 1.5MPa。这一分布规律与\par

\par
DN89、DN65 等单级阀门的压损特征一致 [ \textit{—}] ，即通径突变处的局部阻力是系统压降的主要来源。\par

与文献 [\cite{fe880629_9ce4_4876_a173_e24a97d827f4_基于数字孪生的管道阀门监测技术研究_谭宗豪}]   中 DN89 切断阀全开 12MPa 的压损相比，本系统单级压降更小， 这是由于多级节流结构逐级分配了压降。上游储罐与主管路压力维持在 320MPa 附近，下游出口侧稳定在 220MPa，与边界条件一致，上游高压段压力均匀性良好，说明储罐容腔起到了良好的压力稳压作用 [\cite{fe880629_9ce4_4876_a173_e24a97d827f4_基于数字孪生的管道阀门监测技术研究_谭宗豪}]。\par

\par
图 3    稳态阶段静压力分布云图 [Pa]\par

\par
\textbf{5.1.2    }\textbf{速度分布}\par

速度分布与压力分布呈对应关系，主管内流速较低，平均约为 8 12m/s；节流孔处通径骤缩，流体加速，最大速度可达 58.6m/s，出现在节流孔出口的射流核心区。与文献 [\cite{fe880629_9ce4_4876_a173_e24a97d827f4_基于数字孪生的管道阀门监测技术研究_谭宗豪}]   中 DN89 切断阀全开工况下 134m/s 的最大速度相比，本系统流速更低，这主要是由于多级节流结构逐级降压，单级压降更小所致。\par

三支并联管路中， ms-1 支路通径最大， 流量最高；ms-3 支路通径最小， 流量最低，流量与通径面积近似呈正比关系，实现了分级节流的设计目标 [\cite{fe880629_9ce4_4876_a173_e24a97d827f4_基于数字孪生的管道阀门监测技术研究_谭宗豪}]。\par

\par
\par
图 4    稳态阶段速度大小分布云图 [m/s]\par

\par
\textbf{5.1.3}\textbf{    }\textbf{温度分布}\par

全场温度分布较为均匀， 整体在 288 300K 范围内波动，节流孔下游仅出现约 2 3K的小幅温降。原佳妮 [\cite{fe880629_9ce4_4876_a173_e24a97d827f4_基于数字孪生的管道阀门监测技术研究_谭宗豪}]  在高压球阀研究中指出， 氮气在高压下的焦耳-汤姆逊效应较弱，短流程节流过程的温度变化不显著， 与本仿真结果吻合。与 DN89 阀门 309K 的最高温度相比 [\cite{fe880629_9ce4_4876_a173_e24a97d827f4_基于数字孪生的管道阀门监测技术研究_谭宗豪}] ，本系统温度水平更低， 进一步验证了逐级降压的效果。温度分布结果验证了绝热边界条件设置的合理性，也说明该工况下节流温降效应对系统性能的影响可忽略。\par

\par
图 5    稳态阶段静温度分布云图 [K]\par

\par
\par
\textbf{5.2    }\textbf{阀门瞬态响应特性}\par

两级阀门按预设时序依次动作，对应时间约为 0.085s 与 0.11s，流动呈现典型的瞬态响应特征。\par

\par
\textbf{5.2.1    }\textbf{质量流量响应规律}\par

阀门开启动作后，对应支路的质量流量呈现\textbf{「骤降—超调—回升」}的典型瞬态特征。这一特征由两方面因素共同导致： 一是阀门开启瞬间， 阀后压力尚未建立， 流体惯性导致流量滞后；二是压力波在管路内的反射与叠加， 造成流量的振荡超调。李华贵等 [\cite{fe880629_9ce4_4876_a173_e24a97d827f4_基于数字孪生的管道阀门监测技术研究_谭宗豪}]  在超高压截止阀的动作特性研究中也观测到类似的流量振荡现象， 其本质是压力波传播与流体惯性的耦合作用。\par

具体来看：\par

•  ms-2 支路约在 0.085s 开始响应，流量快速下降后逐步回升，是最先动作的支路；\par

•  ms-1 支路约在 0.11s 开始响应，响应时序与阀门动作设定一致，流量幅值最大；\par

•  ms-3 支路流量始终保持小幅波动，无显著突变，受阀门动作影响最小。\par

三级支路的流量分级特征明显，ms-1 支路流量约为 ms-3 支路的 20 倍，与通径设计比例匹配。上海电气 [\cite{fe880629_9ce4_4876_a173_e24a97d827f4_基于数字孪生的管道阀门监测技术研究_谭宗豪}]  在汇管段仿真中也得到了类似的支路流量分配规律。\par

\textit{·}10\textit{—}2\par

\par
\begin{tabular}{|c|c|c|c|}\n\hline\n &  & Valve action      &  \ \hline\n\end{tabular}
0\textit{.}14     0\textit{.}14      0\textit{.}14     0\textit{.}15      0\textit{.}15     0\textit{.}15      0\textit{.}15      0\textit{.}15     0\textit{.}16      0\textit{.}16     0\textit{.}16 Flow time [s]\par

图 6    各支路质量流量随时间变化曲线\par

\par
\textbf{5.2.2    }\textbf{压力响应特性}\par

阀门动作过程中， 出口压力仅出现小幅波动， 整体维持在 220MPa 附近。三个出口的压力最大波动幅度不超过 50Pa，相对波动量不足百万分之三。梁连金 [\cite{fe880629_9ce4_4876_a173_e24a97d827f4_基于数字孪生的管道阀门监测技术研究_谭宗豪}]  在高压阀门研究中指出， 容腔结构对压力波动具有显著的衰减作用， 本研究结果进一步验证了多级节流系统的中间容腔具有良好的压力缓冲效果， 上游的瞬态扰动经过容腔的储能与耗散后，对下游出口的影响极小。该特性保证了下游负载的压力稳定性， 是多级节流系统的核心优势之一 [\cite{fe880629_9ce4_4876_a173_e24a97d827f4_基于数字孪生的管道阀门监测技术研究_谭宗豪}]。\par

\par
\par
\par
图 7    出口侧监测点压力随时间变化曲线\par

\par
\par
\textbf{6    }\textbf{数值发散问题深入分析}\par

强瞬态阶段的数值发散是本研究遇到的核心问题，结合多轮试算结果与相关理论，从三个层面揭示发散机制：\par

\par
\textbf{6.1    }\textbf{湍流方程耦合发散}\par

k-epsilon 系列湍流模型为半经验模型， 其湍动能产生项与平均速度梯度成正比。阀门动作瞬间产生极强的速度梯度与压力梯度，导致湍动能产生项出现非物理的指数增长，而耗散项无法及时抵消，最终造成 k 与  方程的数值溢出。谭宗豪 [\cite{fe880629_9ce4_4876_a173_e24a97d827f4_基于数字孪生的管道阀门监测技术研究_谭宗豪}]  在阀门泄漏故障仿真中也发现， 强梯度区域的湍动能过度产生是导致计算发散的重要诱因， 这是半经验湍流模型在强非定常强梯度工况下的固有缺陷。原佳妮 [\cite{fe880629_9ce4_4876_a173_e24a97d827f4_基于数字孪生的管道阀门监测技术研究_谭宗豪}]   的研究也表明，瞬态冲击过程中湍流方程的非物理增长是高压阀门仿真中的共性问题。\par

\par
\textbf{6.2    }\textbf{显式求解器的} \textbf{CFL}\textbf{ }\textbf{约束}\par

密度基显式求解器满足 CFL 数 1  的稳定条件， 时间步长必须小于「最小网格尺寸/当地声速」。节流口处网格尺寸小、超高压下当地声速高， 0.001s 的参考步长超出稳定极限约 3 个数量级，直接触发浮点异常。这一现象符合显式时间推进格式的基本稳定理论 [\cite{fe880629_9ce4_4876_a173_e24a97d827f4_基于数字孪生的管道阀门监测技术研究_谭宗豪}] ，说明超高压细网格工况下，显式求解器的计算效率会受到极大限制。\par

\par
\textbf{6.3    }\textbf{几何局部梯度影响}\par

阀芯、节流转角处存在强逆压梯度与流动分离， 是数值误差累积的主要区域。梁连金 [\cite{fe880629_9ce4_4876_a173_e24a97d827f4_基于数字孪生的管道阀门监测技术研究_谭宗豪}]  在阀门强度分析中指出，结构不连续处易产生应力集中与误差放大；流场中同理，\par

\par
大梯度区域的离散误差会随时间步逐步传递放大， 最终触发全局发散。低质量网格与大梯度流动的耦合作用，是数值稳定性下降的重要诱因。\par

\par
\textbf{7    }\textbf{结论与展望}\par

\par
\textbf{7.1    }\textbf{主要结论}\par

1. 超高压多级节流系统的稳态压降主要集中在节流口位置， 约占总压降的 85\\%；节流口处流速显著升高，最大速度达 58.6m/s，主管压力损失小，压损分布规律与 DN89、 DN65 等单级阀门的工程数据一致 [ \textit{—}] ；2.  阀门动作过程中， 支路流量呈现「骤降—超调—回升」的典型瞬态特征， 两级阀门响应时序与设定一致，响应时间约 0.03s 量级，流量特征与压力波传播、流体惯性的耦合作用相关， 与已有文献结论一致 [\cite{fe880629_9ce4_4876_a173_e24a97d827f4_基于数字孪生的管道阀门监测技术研究_谭宗豪}] ；3.  系统容腔对下游压力具有显著的缓冲作用， 阀门动作过程中出口压力波动小于 50Pa，相对波动量不足百万分之三， 下游压力稳定性优异， 验证了容腔缓冲结构的有效性 [\cite{fe880629_9ce4_4876_a173_e24a97d827f4_基于数字孪生的管道阀门监测技术研究_谭宗豪}] ；4. 强瞬态工况下，k-epsilon 湍流模型的数值稳定性差， 易出现湍动能溢出； 显式密度基求解器对局部细网格的时间步长要求严苛，起步易发散；层流模型可获得稳定的定性结果。\par

\par
\textbf{7.2    }\textbf{后续研究展望}\par

1.  采用 k-omega SST 湍流模型替代 k-epsilon 模型，利用其强逆压梯度下的稳定性优势，改善阀门动作阶段的收敛性。谭宗豪 [\cite{fe880629_9ce4_4876_a173_e24a97d827f4_基于数字孪生的管道阀门监测技术研究_谭宗豪}]   的研究表明该模型对强梯度流动的适应性更优；2.  采用变时间步长策略，稳态阶段用大步长提速， 阀门动作阶段自动加密步长，兼顾效率与稳定性；3.  对节流口、阀芯区域进行网格平滑优化，降低单元纵横比，消除低质量网格。参考曲率尺寸函数方法 [\cite{fe880629_9ce4_4876_a173_e24a97d827f4_基于数字孪生的管道阀门监测技术研究_谭宗豪}]  提升网格质量；4.  结合孔口流量公式与气体动力学理论， 对仿真结果进行半理论校验， 进一步提升结论可靠性。同时可参考李华贵等 [\cite{fe880629_9ce4_4876_a173_e24a97d827f4_基于数字孪生的管道阀门监测技术研究_谭宗豪}]  的结构优化思路，开展阀门结构与流场的耦合优化。\par

\par
\textbf{8    }\textbf{参考文献}\par

[\cite{fe880629_9ce4_4876_a173_e24a97d827f4_基于数字孪生的管道阀门监测技术研究_谭宗豪}]  DN89 切断阀压损稳态仿真阶段性报告 [R].  上海:  上海自动化仪表有限公司, 2025.\par

[\cite{fe880629_9ce4_4876_a173_e24a97d827f4_基于数字孪生的管道阀门监测技术研究_谭宗豪}] DN65 阀门稳态压损仿真阶段性报告 [R]. 上海: 上海自动化仪表有限公司, 2025.\par

[\cite{fe880629_9ce4_4876_a173_e24a97d827f4_基于数字孪生的管道阀门监测技术研究_谭宗豪}]  李华贵,  罗建康,  项良海,  等.   超高压大口径气动快速截止阀的研制 [J].  阀门, 2025(7):  713-719.\par

[\cite{fe880629_9ce4_4876_a173_e24a97d827f4_基于数字孪生的管道阀门监测技术研究_谭宗豪}] 上海电气. 超高压阀门系统仿真汇报 [R]. 上海: 上海自动化仪表有限公司, 2025.\par

[\cite{fe880629_9ce4_4876_a173_e24a97d827f4_基于数字孪生的管道阀门监测技术研究_谭宗豪}] 原佳妮.  高温高压气动球阀的结构优化及响应特性研究 [D]. 西安:  西安工业大学, 2025.\par

\par
[\cite{fe880629_9ce4_4876_a173_e24a97d827f4_基于数字孪生的管道阀门监测技术研究_谭宗豪}] 梁连金.  高压大口径双向密封程控提升杆球阀的研制 [J]. 石油化工设计, 2025, 42(2):  15-21.\par

[\cite{fe880629_9ce4_4876_a173_e24a97d827f4_基于数字孪生的管道阀门监测技术研究_谭宗豪}] 谭宗豪. 基于数字孪生的管道阀门监测技术研究 [D]. 东莞: 东莞理工学院, 2025.\par


\printbibliography
\end{document}
